% !TeX root = ../tcc.tex
\chapter{Introdução}

Este capítulo apresenta o contexto da aplicação de modelos de visão
computacional à identificação de manifestações dermatológicas associadas
à hanseníase. São delimitados o problema de pesquisa, a questão
norteadora e os objetivos do trabalho, considerando a avaliação de
desempenho na identificação de padrões associados à hanseníase em
imagens de pele, no contexto de apoio à triagem na rede pública de saúde
do Brasil.

\section{Contexto}\label{contexto}

A hanseníase integra o grupo das doenças tropicais negligenciadas e
constitui um desafio para a identificação oportuna de casos e o
acompanhamento das pessoas afetadas. Em sua revisão da literatura,
\citeonline{fernandes2024} destacam que o diagnóstico se baseia
principalmente na avaliação clínica e que a semelhança com outras
condições dermatológicas pode dificultar o reconhecimento da doença.
Nesse contexto, a investigação de recursos computacionais de apoio à
análise de manifestações cutâneas pode contribuir para ampliar as
possibilidades de suporte aos profissionais de saúde, desde que sua
avaliação considere as características da tarefa e da população de
interesse.

A análise automatizada de imagens dermatológicas tem sido investigada
por meio de técnicas de aprendizado de máquina e aprendizado profundo.
\citeonline{fernandes2024} identificaram 21 estudos sobre inteligência
artificial no apoio ao diagnóstico da hanseníase e observaram que a
maior parte utilizava imagens de lesões cutâneas e tarefas de
classificação multiclasse. A inclusão da hanseníase em sistemas voltados
a diferentes doenças indica que os resultados dessas propostas devem ser
interpretados tanto em relação ao conjunto de classes quanto ao
reconhecimento específico da doença.

Entre os estudos experimentais, \citeonline{jaikishore2021}
desenvolveram uma abordagem baseada em MobileNetV2 para classificar
imagens de eczema, sarampo, hanseníase e pele saudável, incorporando o
modelo a um protótipo Android. \citeonline{yotsu2023}, por sua vez,
avaliaram ResNet-50 e VGG-16 na classificação de cinco doenças tropicais
negligenciadas com manifestações cutâneas. Em outra linha de
investigação, \citeonline{eita2025} propuseram um sistema baseado em YOLOv8
para localizar e classificar lesões de diferentes condições, incluindo
hanseníase, com uso de mapas de saliência para apoiar a interpretação
das predições. Esses trabalhos ilustram abordagens distintas de análise
de imagens, com diferentes requisitos de dados, saídas e procedimentos
de avaliação.

O campo também compreende tarefas relacionadas ao acompanhamento de
manifestações e complicações da doença. \citeonline{sanchez2024}
apresentaram o CO2Wounds-V2, um conjunto de fotografias de feridas
crônicas de pacientes com hanseníase, acompanhado de experimentos de
segmentação semântica. Nesse caso, a tarefa consiste em delimitar
regiões de ferida nas imagens, e não em determinar se o paciente
apresenta hanseníase. Essa distinção é necessária para organizar a
literatura e evitar que medidas de segmentação sejam interpretadas como
evidências de desempenho diagnóstico.

A caracterização das populações representadas nos dados também integra a
avaliação dos modelos. \citeonline{yotsu2023} utilizaram fotografias
coletadas em Costa do Marfim e Gana, descrevendo a população como
composta por pessoas de pele escura, com fototipos de Fitzpatrick IV ou
superiores. O estudo oferece evidências diretamente relacionadas à
análise de imagens de pele escura, mas não apresenta comparação de
desempenho entre os diferentes fototipos. \citeonline{eita2025} também
reconhecem limitações na diversidade dos dados e apontam a ampliação da
representação de pessoas com tons de pele mais escuros como uma direção
de trabalho futuro.

Esses estudos motivam a investigação de como o desempenho dos modelos é
definido, mensurado e interpretado. Para o presente trabalho, a
avaliação compreende os resultados preditivos relativos à tarefa, os
procedimentos empregados para obtê-los e as evidências disponíveis sobre
sua aplicação a novos dados e a diferentes grupos de tom de pele. A
utilização futura em apoio à triagem na rede pública de saúde do Brasil
constitui o contexto de interesse da pesquisa, sem que os resultados
computacionais sejam, por si só, considerados demonstração de
efetividade clínica.

\section{Problema}\label{problema}

A existência de modelos capazes de reconhecer padrões em imagens não
estabelece, isoladamente, sua adequação à identificação de manifestações
de hanseníase em uma população específica. Os estudos analisados
apresentam diferenças nas tarefas, nas classes de interesse, na
composição dos conjuntos de dados e nas métricas utilizadas. Classificar
uma imagem entre doenças, localizar uma lesão e segmentar uma ferida
produzem resultados com significados distintos. Consequentemente, uma
comparação baseada apenas nos maiores valores publicados pode
desconsiderar diferenças essenciais entre os experimentos \cite{yotsu2023,sanchez2024,eita2025}.

Além disso, medidas agregadas podem não representar adequadamente o
desempenho referente à hanseníase. No estudo de \citeonline{yotsu2023}, a
Task 2 apresentou acurácias globais de 77\% e 80\% nos dois cenários
avaliados. Entretanto, a leitura da diagonal correspondente à hanseníase
nas matrizes de confusão indica proporções de classificação correta de
12\% e 13\%, respectivamente. Como as linhas representam a classe real,
esses valores podem ser interpretados como sensibilidade por imagem para
a classe hanseníase. O contraste evidencia a importância de examinar
resultados específicos por doença, em vez de atribuir à hanseníase o
desempenho global do sistema.

A relação entre os resultados obtidos e os tons de pele constitui outra
dimensão do problema. A presença de imagens de pele escura em um
conjunto de dados não demonstra que o modelo apresente comportamento
equivalente entre diferentes tons. No estudo de \citeonline{yotsu2023}, a
população é caracterizada por um intervalo de fototipos, mas não são
fornecidos resultados separados por esses grupos. Já \citeonline{eita2025}
reconhecem limitações de representatividade e a necessidade de ampliar a
diversidade dos dados. Essas observações sustentam a necessidade de
examinar a caracterização populacional e a avaliação por grupos, sem
permitir concluir, a partir dessas evidências isoladas, que todos os
modelos tenham desempenho inferior em pele negra.

Também permanece necessário distinguir a avaliação computacional da
avaliação em contexto assistencial. O desenvolvimento de um aplicativo
ou a obtenção de resultados em conjuntos de imagens não comprova sua
contribuição para o atendimento em serviços de saúde. Enquanto
\citeonline{jaikishore2021} apresentam um protótipo de
aplicação móvel, \citeonline{eita2025} situam os estudos prospectivos com
dermatologistas entre as atividades futuras. Assim, a relevância dessas
propostas para a rede pública exige considerar o alcance e os limites
dos procedimentos de validação relatados.

Diante desse cenário, o problema investigado consiste em compreender
quais evidências e procedimentos permitem avaliar o desempenho de
modelos de visão computacional na identificação de manifestações
dermatológicas associadas à hanseníase, no contexto de apoio à triagem
na rede pública de saúde do Brasil, considerando a tarefa realizada, as
características dos dados e os limites de generalização dos resultados.
A comparação entre tons de pele será examinada quando houver informações
disponíveis, como uma dimensão complementar da avaliação. A revisão em
andamento buscará sistematizar essas informações e identificar os
aspectos que precisam ser considerados no planejamento de avaliações
futuras.

\section{Questão de pesquisa}\label{gqm-adaptado-e-questuxe3o-de-pesquisa}

\subsection{GQM adaptado}

O objetivo de análise foi organizado por meio de um GQM adaptado, que
explicita o objeto, o propósito, o foco, o ponto de vista e o contexto
da investigação. O quadro orienta a definição da questão de pesquisa e
das medidas a serem extraídas dos estudos.

{\def\LTcaptype{none} % do not increment counter
\begin{longtable}[]{@{}
  >{\raggedright\arraybackslash}p{(\linewidth - 2\tabcolsep) * \real{0.5000}}
  >{\raggedright\arraybackslash}p{(\linewidth - 2\tabcolsep) * \real{0.5000}}@{}}
\toprule\noalign{}
\begin{minipage}[b]{\linewidth}\raggedright
Característica
\end{minipage} & \begin{minipage}[b]{\linewidth}\raggedright
Valor
\end{minipage} \\
\midrule\noalign{}
\endhead
\bottomrule\noalign{}
\endlastfoot
\textbf{Analisar} & Modelos de visão computacional aplicados à
identificação de manifestações dermatológicas. \\
\textbf{Com o propósito de} & Avaliar o desempenho dos modelos aplicados
à identificação de padrões em imagens de pele. \\
\textbf{Com foco em} & Identificação de características patológicas de
hanseníase. \\
\textbf{Do ponto de vista de} & Pesquisadores das áreas de Engenharia de
Software e da área da saúde. \\
\textbf{No contexto de} & Apoio à triagem de hanseníase na rede pública
de saúde do Brasil. \\
\end{longtable}
}

\subsection{Questão norteadora}

A investigação é orientada pela seguinte questão:

\begin{quote}
Como avaliar o desempenho de modelos de visão computacional aplicados à
identificação de padrões associados à hanseníase em imagens de pele, no
contexto da rede pública de saúde do Brasil?
\end{quote}

Nesta etapa, a questão será abordada por meio da identificação e análise
das evidências apresentadas na literatura. Serão examinadas as tarefas
realizadas, os modelos empregados, as características dos conjuntos de
dados, os procedimentos de avaliação, as métricas e os resultados, bem
como as informações disponíveis sobre tons de pele e generalização. A
ausência de avaliação por grupo será registrada como uma limitação das
evidências disponíveis, sem ser tomada como demonstração de equidade ou
de desigualdade de desempenho.

\section{Objetivos}\label{objetivos}

\subsection{Objetivo geral}\label{objetivo-geral}

Analisar a avaliação de desempenho de modelos de visão computacional
aplicados à identificação de padrões associados à hanseníase em imagens
de pele, a partir da síntese das evidências da literatura, visando
fundamentar seu uso em apoio à triagem na rede pública de saúde do
Brasil.

\subsection{Objetivos específicos}\label{objetivos-especuxedficos}

\begin{enumerate}
\def\labelenumi{\arabic{enumi}.}
\tightlist
\item
  Caracterizar as tarefas, os modelos e as abordagens de visão
  computacional empregados na análise de manifestações dermatológicas
  associadas à hanseníase.
\item
  Identificar os conjuntos de dados utilizados, suas formas de aquisição
  e anotação e as informações disponíveis sobre as populações e os tons
  de pele representados.
\item
  Sistematizar os procedimentos de avaliação, as métricas e os
  resultados reportados, distinguindo medidas globais, resultados
  específicos de hanseníase e unidades de análise.
\item
  Examinar as evidências de generalização e de avaliação do desempenho
  entre grupos de tom de pele, quando disponíveis.
\item
  Identificar limitações e lacunas que fundamentem o planejamento de uma
  avaliação futura de modelos no contexto de apoio à triagem na rede
  pública de saúde do Brasil.
\end{enumerate}


\section{Organização do trabalho}

Este trabalho está organizado em três capítulos. O primeiro apresenta o
contexto, o problema de pesquisa, a questão norteadora e os objetivos. O
segundo reúne os fundamentos sobre pesquisa empírica, hanseníase, visão
computacional e avaliação de modelos. O terceiro descreve o planejamento,
a busca e a seleção dos estudos da revisão de literatura, bem como os
procedimentos previstos para a extração e a síntese das evidências. Ao
final, são apresentadas as referências e os apêndices.

\section{Cronograma e fluxo de atividades}

O trabalho está organizado em etapas que compreendem o planejamento da
pesquisa, a revisão da literatura e a análise das evidências obtidas. O
fluxo apresentado a seguir descreve a sequência das atividades; as datas,
a duração de cada etapa e eventuais atividades experimentais posteriores
devem ser definidas com a orientação.

\subsection{Atividades da primeira etapa}

Na primeira etapa, são estabelecidos o problema, a questão de pesquisa e os
objetivos; estruturado o protocolo da revisão; e realizada a identificação e
a triagem inicial dos estudos. Em seguida, estão previstas a leitura integral
dos registros selecionados, a confirmação de sua elegibilidade, a extração
das informações e a síntese dos resultados. O estado da triagem e as
pendências são apresentados no Capítulo~\ref{cap:revisao-literatura}.

\subsection{Atividades da segunda etapa}

A continuidade do trabalho dependerá dos resultados da revisão e das
definições acordadas com a orientação. Entre as atividades a planejar estão
o aprofundamento da análise das evidências e a definição de procedimentos
para uma avaliação futura de modelos no contexto de apoio à triagem. A
realização de experimentos, a disponibilidade de dados e os prazos dessa
etapa ainda precisam ser estabelecidos.
